% titel: Blatt 1, A2: Kinematik -- Protonen-Beschleuniger
% skript: Kap. Kinematik, Bewegung in einer Dimension

% Stueckweise Beschleunigung durch Integration zu v(t) und x(t).

Bei bekannter Beschleunigung $a(t)$ einer eindimensionalen Bewegung
gewinnt man Geschwindigkeit und Ort durch zweimaliges Integrieren.

\begin{werkzeugkasten}[Werkzeuge]
Aus $a(t)=a$ folgt durch Integration mit den Anfangsbedingungen
$v(0)=0$, $x(0)=0$:
\begin{equation}
  v(t)=v(0)+\int_0^t a(t')\,\mathrm dt',\qquad
  x(t)=x(0)+\int_0^t v(t')\,\mathrm dt'\;.
\end{equation}

\textbf{Stueckweise Funktionen} loest man Intervall fuer Intervall: die
Werte am rechten Rand eines Intervalls dienen als Anfangswerte des
naechsten. Der Anschluss muss stetig sein.

\textbf{Schaetzungen:} Bei einer Spitze, die linear auf $a_0$ ansteigt und
linear abfaellt, ist die Flaeche $\int a\,\mathrm dt$ (die
Geschwindigkeitsaenderung) die halbe Rechteckflaeche.
\end{werkzeugkasten}

\begin{werkzeugkasten}[Vorgehen]
\begin{itemize}
  \item Die drei Intervalle $0\le t<T$, $T\le t<2T$, $t\ge2T$ getrennt
        aufschreiben.
  \item Erst $a(t)$ in jedem Intervall als Geradengleichung angeben, dann
        integrieren.
  \item Die Endgeschwindigkeit liest man nach Abklingen der Beschleunigung
        direkt ab, da sie ab $t=2T$ konstant bleibt.
\end{itemize}
\end{werkzeugkasten}

\begin{achtung}[Typische Fallen]
\begin{itemize}
  \item Integrationskonstanten nicht vergessen; sie werden durch die
        Anschlussbedingungen festgelegt.
  \item $x(t)$ erhaelt man durch Integrieren von $v(t)$, nicht von $a(t)$.
  \item Die Geschwindigkeit darf keine Spruenge aufweisen, auch wenn die
        Beschleunigung sprunghaft ist.
\end{itemize}
\end{achtung}
